\section{Fourier inversion and exactness}\label{sec:fourier}

For a subset $A\subseteq\mathcal E$, give $A$ the Bernoulli weight
\[
w(A)=\theta^{|A|}(1-\theta)^{M-|A|}.
\tag{3.1}
\]
Let $P$ be the total weight of exact representations of $1/b$. We first work in the finite group $\Z/L\Z$.

\subsection{The target coefficient}

Let $P^{\mathrm{quot}}$ be the probability of the congruence
\[
\sum_{e\in\mathcal E}\xi_e\frac{L}{e}
\equiv\frac{L}{b}\pmod L.
\tag{3.2}
\]
Character orthogonality gives
\[
L P^{\mathrm{quot}}
=
\sum_{h\bmod L}F(h),
\tag{3.3}
\]
where
\[
F(h)
=
e(-h/b)
\prod_{e\in\mathcal E}
\bigl((1-\theta)+\theta e(h/e)\bigr),
\qquad
e(t)=e^{2\pi it}.
\tag{3.4}
\]

We will use the following modulus estimate repeatedly.

\begin{lemma}[Bernoulli modulus bound]\label{lem:bernoulli-modulus}
There is $\kappa_b>0$ such that for every real $u$,
\[
\left|(1-\theta)+\theta e(u)\right|
\le
\exp(-\kappa_b\|u\|^2),
\tag{3.5}
\]
where $\|u\|$ denotes the distance from $u$ to the nearest integer.
\end{lemma}

\begin{proof}
We have
\[
\left|(1-\theta)+\theta e(u)\right|^2
=
1-4\theta(1-\theta)\sin^2(\pi u).
\]
The parameter $\theta$ lies in a fixed compact subinterval of $(0,1)$, while
$\sin(\pi x)\ge2x$ for $0\le x\le1/2$. Thus
\[
4\theta(1-\theta)\sin^2(\pi u)
\ge c_b\|u\|^2.
\]
Using $1-v\le e^{-v}$ and adjusting the constant after taking square roots proves \eqref{3.5}.
\end{proof}

In particular,
\[
|F(h)|
\le
\exp\!\left(-\kappa_b
\sum_{e\in\mathcal E}\|h/e\|^2\right).
\tag{3.6}
\]
The problem is therefore geometric: we must understand which CRT configurations make a large collection of the phases $h/e$ simultaneously close to integers.

\subsection{CRT coordinates and the factor partition}

Since $L$ is squarefree, CRT identifies a frequency $h\bmod L$ with one residue coordinate modulo every prime dividing $L$. The denominator family naturally separates these coordinates.

The primes in $\mathcal B$ are the \emph{anchor coordinates}. Denominators $qq'$ with $q,q'\in\mathcal B$ measure their internal coherence. Once the anchor is coherent, denominators $rq$ with $q\in\mathcal B$ and
\[
r\in(\mathcal P\setminus\mathcal B)\cup S_b
\]
observe one remaining row against the common anchor. Denominators involving no anchor prime are deliberately left unused by the decoder and will later provide extra damping.

Every denominator in $\mathcal E$ belongs to exactly one of the classes
\begin{enumerate}[label=\textup{(\roman*)}]
\item $qq'$ with $q,q'\in\mathcal B$: anchor--anchor factors;
\item $rq$ with $q\in\mathcal B$, $r\in\mathcal P\setminus\mathcal B$: lower rows;
\item $rq$ with $q\in\mathcal B$, $r\in S_b$: target rows;
\item $pp'$ with $p,p'\in\mathcal P\setminus\mathcal B$: retained lower--lower factors.
\end{enumerate}
This exact partition will be load-bearing in Section \ref{sec:rows}: classes (i)--(iii) recover the coherent CRT state, while class (iv) suppresses coherent labels outside the Gaussian range.

\subsection{Genuine integer frequencies}

Every integer $m$ determines the CRT assignment obtained by reducing $m$ modulo each prime dividing $L$. When $|m|<e/2$ for every $e\in\mathcal E$,
\[
\left\|\frac{m}{e}\right\|=\frac{|m|}{e}.
\tag{3.7}
\]
Moreover the linear phase in the logarithm of \eqref{3.4} cancels exactly because
\[
\theta\sum_{e\in\mathcal E}\frac1e=\frac1b.
\tag{3.8}
\]
This produces the Gaussian main term in Section \ref{sec:gaussian}. Before the Taylor expansion can control the full Fourier sum, however, we must show that every sufficiently low-energy CRT configuration is not merely close to an integer assignment but is exactly one. Section \ref{sec:anchor} proves this for the anchor block and Section \ref{sec:rows} for the remaining rows.

\subsection{No-wrap as a separate exactness step}

The Fourier calculation takes place in $\Z/L\Z$, whereas the theorem is an equality in $\Q$. We keep these two assertions logically separate.

\begin{lemma}[No-wrap]\label{lem:no-wrap}
If $A\subseteq\mathcal E$ satisfies
\[
\sum_{e\in A}\frac{L}{e}
\equiv\frac{L}{b}\pmod L,
\tag{3.9}
\]
then
\[
\sum_{e\in A}\frac1e=\frac1b.
\tag{3.10}
\]
\end{lemma}

\begin{proof}
After division by $L$,
\[
D=
\sum_{e\in A}\frac1e-\frac1b
\]
is an integer. On the other hand,
\[
0\le\sum_{e\in A}\frac1e\le\Lambda<1,
\qquad
0<\frac1b<1.
\]
Thus $-1<D<1$, so $D=0$.
\end{proof}

Consequently
\[
P^{\mathrm{quot}}=P.
\tag{3.11}
\]
We nevertheless retain the quotient notation during the Fourier estimates to emphasize that positivity in the finite group and exactness in the fundamental interval are different inputs.

The same observation will apply to moving targets. If
\[
t_j=\frac1b+\frac{j}{L}
\]
lies in $[0,\Lambda]$, then a subset sum congruent to $t_j$ modulo one is equal to $t_j$. For every fixed standardized radius $U$, this condition holds for all $|j|\le UL\sigma$ once $X$ is sufficiently large, because $\sigma\to0$ and $1/b$ lies strictly inside the limiting load interval.
